\chapter{2024年新高考II卷}

\section{单选题}

\begin{problem}
已知 \(z=-1-\i\),则 \(|z|=\)(\quad)

\choices
    {\(0\)}
    {\(1\)}
    {\(\sqrt2\)}
    {\(2\)}
\end{problem}

\begin{problem}
已知命题 \(p:\forall x\in\R\),\(|x+1|>1\);命题 \(q:\exists x>0\),\(x^3=x\),则(\quad)

\choices
    {\(p\) 和 \(q\) 都是真命题}
    {\(\neg p\) 和 \(q\) 都是真命题}
    {\(p\) 和 \(\neg q\) 都是真命题}
    {\(\neg p\) 和 \(\neg q\) 都是真命题}
\end{problem}

\begin{problem}
已知向量 \(\bs a\),\(\bs b\) 满足 \(|\bs a|=1\),\(|\bs a+2\bs b|=2\),且 \((\bs b-2\bs a)\perp \bs b\),则 \(|\bs b|=\)(\quad)

\choices
    {\(\frac12\)}
    {\(\frac{\sqrt2}{2}\)}
    {\(\frac{\sqrt3}{2}\)}
    {\(1\)}
\end{problem}

\begin{problem}
某农业研究部门在面积相等的 \(100\text{ 块}\)稻田上种植一种新型水稻,得到各块稻田的亩产量(单位:\(\mathrm{kg}\))并整理得如下表. 根据表中数据,结论中正确的是(\quad)

\begin{center}
\begin{tabular}{|c|c|c|c|}
\hline
亩产量 & \([900,950)\) & \([950,1000)\) & \([1000,1050)\) \\
\hline
频数 & \(6\) & \(12\) & \(18\) \\
\hline
亩产量 & \([1050,1100)\) & \([1100,1150)\) & \([1150,1200)\) \\
\hline
频数 & \(30\) & \(24\) & \(10\) \\
\hline
\end{tabular}
\end{center}

\choices
    {\(100\text{ 块}\)稻田亩产量的中位数小于 \(1050 \mathrm{~kg}\)}
    {\(100\text{ 块}\)稻田中亩产量低于 \(1100 \mathrm{~kg}\) 的稻田所占比例超过 \(80\%\)}
    {\(100\text{ 块}\)稻田亩产量的极差介于 \(200 \mathrm{~kg}\) 至 \(300 \mathrm{~kg}\) 之间}
    {\(100\text{ 块}\)稻田亩产量的平均值介于 \(900 \mathrm{~kg}\) 至 \(1000 \mathrm{~kg}\) 之间}
\end{problem}

\begin{problem}
已知曲线 \(C:x^2+y^2=16\)(\(y>0\)),从 \(C\) 上任意一点 \(P\) 向 \(x\) 轴作垂线段 \(PP'\),\(P'\) 为垂足,则线段 \(PP'\) 的中点 \(M\) 的轨迹方程为(\quad)

\choices
    {\(\dfrac{x^2}{16}+\dfrac{y^2}{4}=1\)(\(y>0\))}
    {\(\dfrac{x^2}{16}+\dfrac{y^2}{8}=1\)(\(y>0\))}
    {\(\dfrac{y^2}{16}+\dfrac{x^2}{4}=1\)(\(y>0\))}
    {\(\dfrac{y^2}{16}+\dfrac{x^2}{8}=1\)(\(y>0\))}
\end{problem}

\begin{problem}
设函数 \(f(x)=a(x+1)^2-1\),\(g(x)=\cos x+2ax\),当 \(x\in(-1,1)\) 时,曲线 \(y=f(x)\) 与 \(y=g(x)\) 恰有一个交点,则 \(a=\)(\quad)

\choices
    {\(-1\)}
    {\(\frac12\)}
    {\(1\)}
    {\(2\)}
\end{problem}

\begin{problem}
已知正三棱台 \(ABC-A_1B_1C_1\) 的体积为 \(\frac{52}{3}\),\(AB=6\),\(A_1B_1=2\),则 \(A_1A\) 与平面 \(ABC\) 所成角的正切值为(\quad)

\choices
    {\(\frac12\)}
    {\(1\)}
    {\(2\)}
    {\(3\)}
\end{problem}

\begin{problem}
设函数 \(f(x)=(x+a)\ln(x+b)\),若 \(f(x)\ge 0\),则 \(a^2+b^2\) 的最小值为(\quad)

\choices
    {\(\frac18\)}
    {\(\frac14\)}
    {\(\frac12\)}
    {\(1\)}
\end{problem}

\section{多选题}

\begin{problem}
对于函数 \(f(x)=\sin 2x\) 和 \(g(x)=\sin\left(2x-\frac{\pi}{4}\right)\),下列说法正确的有(\quad)

\choices
    {\(f(x)\) 与 \(g(x)\) 有相同的零点}
    {\(f(x)\) 与 \(g(x)\) 有相同的最大值}
    {\(f(x)\) 与 \(g(x)\) 有相同的最小正周期}
    {\(f(x)\) 与 \(g(x)\) 的图像有相同的对称轴}
\end{problem}

\begin{problem}
抛物线 \(C:y^2=4x\) 的准线为 \(l\),\(P\) 为 \(C\) 上的动点,过 \(P\) 作圆 \(\odot A:x^2+(y-4)^2=1\) 的一条切线,\(Q\) 为切点;过 \(P\) 作 \(l\) 的垂线,垂足为 \(B\),则(\quad)

\choices
    {\(l\) 与 \(\odot A\) 相切}
    {当 \(P\),\(A\),\(B\) 三点共线时,\(|PQ|=\sqrt{15}\)}
    {当 \(|PB|=2\) 时,\(PA\perp AB\)}
    {满足 \(|PA|=|PB|\) 的点 \(P\) 有且仅有 \(2\) 个}
\end{problem}

\begin{problem}
设函数 \(f(x)=2x^3-3ax^2+1\),则(\quad)

\choices
    {当 \(a>1\) 时,\(f(x)\) 有三个零点}
    {当 \(a<0\) 时,\(x=0\) 是 \(f(x)\) 的极大值点}
    {存在 \(a\),\(b\),使得 \(x=b\) 为曲线 \(y=f(x)\) 的对称轴}
    {存在 \(a\),使得点 \((1,f(1))\) 为曲线 \(y=f(x)\) 的对称中心}
\end{problem}

\section{填空题}

\begin{problem}
记 \(S_n\) 为等差数列 \(\{a_n\}\) 的前 \(n\) 项和,若 \(a_3+a_4=7\),\(3a_2+a_5=5\),则 \(S_{10}=\)\fillinblank{}.
\end{problem}

\begin{problem}
已知 \(\alpha\) 为第一象限角,\(\beta\) 为第三象限角,\(\tan\alpha+\tan\beta=4\),\(\tan\alpha\tan\beta=\sqrt2+1\),则 \(\sin(\alpha+\beta)=\)\fillinblank{}.
\end{problem}

\begin{problem}
在如图的 \(4\times 4\) 方格表中选 4 个方格,要求每行和每列均恰有一个方格被选中,则共有\fillinblank{} 种选法;在所有符合上述要求的选法中,选中方格中的 4 个数之和的最大值是\fillinblank{}.

\begin{center}
\begin{tabular}{|c|c|c|c|}
\hline
11 & 21 & 31 & 40 \\
\hline
12 & 22 & 33 & 42 \\
\hline
13 & 22 & 33 & 43 \\
\hline
15 & 24 & 34 & 44 \\
\hline
\end{tabular}
\end{center}
\end{problem}

\section{解答题}

\begin{problem}
记 \(\triangle ABC\) 的内角 \(A\),\(B\),\(C\) 的对边分别为 \(a\),\(b\),\(c\),已知 \(\sin A+\sqrt3\cos A=2\).
\begin{enumerate}
\item 求 \(A\);
\item 若 \(a=2\),\(\sqrt2 b\sin C=c\sin 2B\),求 \(\triangle ABC\) 的周长.
\end{enumerate}
\end{problem}

\begin{problem}
已知函数 \(f(x)=\e^x-ax-a^3\).
\begin{enumerate}
\item 当 \(a=1\) 时,求曲线 \(y=f(x)\) 在点 \((1,f(1))\) 处的切线方程;
\item 若 \(f(x)\) 有极小值,且极小值小于 0,求 \(a\) 的取值范围.
\end{enumerate}
\end{problem}

\begin{problem}
如图,平面四边形 \(ABCD\) 中,\(AB=8\),\(CD=3\),\(AD=5\sqrt3\),\(\angle ADC=90^\circ\),\(\angle BAD=30^\circ\),点 \(E\),\(F\) 满足 \(\overrightarrow{AE}=\frac{2}{5}\overrightarrow{AD}\),\(\overrightarrow{AF}=\frac{1}{2}\overrightarrow{AB}\),将 \(\triangle AEF\) 沿 \(EF\) 对折至 \(\triangle PEF\),使得 \(PC=4\sqrt3\).
\begin{enumerate}
\item 证明:\(EF\perp PD\);
\item 求面 \(PCD\) 与面 \(PBF\) 所成二面角的正弦值.
\end{enumerate}

\bitmapfigure[max width=4.6cm]{img/2024/new_gaokao_paper_2/q17_fig1.png}
\end{problem}

\begin{problem}
某投篮比赛分为两个阶段,每个参赛队由两名队员组成,比赛具体规则如下:第一阶段由参赛队中一名队员投篮 \(3\text{ 次}\),若 \(3\text{ 次}\)都未投中,则该队被淘汰,比赛成绩为 \(0\text{ 分}\);若至少投中一次,则该队进入第二阶段,由该队的另一名队员投篮 \(3\text{ 次}\),每次投中得 \(5\text{ 分}\),未投中得 \(0\text{ 分}\),该队的比赛成绩为第二阶段的得分总和. 某参赛队由甲,乙两名队员组成,设甲每次投中的概率为 \(p\),乙每次投中的概率为 \(q\),各次投中与否相互独立.
\begin{enumerate}
\item 若 \(p=0.4\),\(q=0.5\),甲参加第一阶段比赛,求甲,乙所在队的比赛成绩不少于 5 分的概率;
\item 假设 \(0<p<q\),
\begin{enumerate}
\item 为使得甲,乙所在队的比赛成绩为 15 分的概率最大,应该由谁参加第一阶段比赛?
\item 为使得甲,乙所在队的比赛成绩的数学期望最大,应该由谁参加第一阶段比赛?
\end{enumerate}
\end{enumerate}
\end{problem}

\begin{problem}
已知双曲线 \(C:x^2-y^2=m\)(\(m>0\)),点 \(P_1(5,4)\) 在 \(C\) 上,\(k\) 为常数,\(0<k<1\). 按照如下方式依次构造点 \(P_n\)(\(n=2\),\(3\),\(\dots\)):过 \(P_{n-1}\) 作斜率为 \(k\) 的直线与 \(C\) 的左支交于点 \(Q_{n-1}\),令 \(P_n\) 为 \(Q_{n-1}\) 关于 \(y\) 轴的对称点,记 \(P_n\) 的坐标为 \((x_n,y_n)\).
\begin{enumerate}
\item 若 \(k=\frac12\),求 \(x_2\),\(y_2\);
\item 证明:数列 \(\{x_n-y_n\}\) 是公比为 \(\frac{1+k}{1-k}\) 的等比数列;
\item 设 \(S_n\) 为 \(\triangle P_nP_{n+1}P_{n+2}\) 的面积,证明:对任意的正整数 \(n\),\(S_n=S_{n+1}\).
\end{enumerate}
\end{problem}

